\documentclass[11pt]{letter}

\usepackage[margin=0.75in]{geometry}
\usepackage{setspace}
\usepackage{hyperref}
\usepackage{xcolor}
\usepackage{enumitem}
\usepackage{parskip}

\hypersetup{
    colorlinks=true,
    urlcolor=black,
    linkcolor=black
}

\setlength{\parindent}{0pt}
\setstretch{1.05}
\date{}
\begin{document}

\begin{letter}{}

\opening{Dear Hiring Manager,}

I saw the job post shared through the CEAS graduate students list and thought I'd be a good fit for it. I am a Master's student in Computer Science at the University of Cincinnati, graduating December 2026, applying for the \textbf{Software Developer Intern} role at \textbf{Global Quality Corp}. The role interests me because it is about working inside a real, evolving codebase rather than a clean-slate project. The construction management RPG in Unreal Engine 5, with its focus on C++ gameplay systems, foundational architecture, and tools, is exactly the kind of long-lived system where careful design and steady ownership matter. I have strong C and C++ and object-oriented fundamentals, I pick up unfamiliar stacks quickly, and I have full-stack JavaScript and React experience that maps to the browser-based 3D VR surgery planning work as well. I would join to learn the project deeply and take on more implementation and technical problem-solving as I go, which is how the role is described.

The clearest example of this from my own work is my time at Siemens Industry Software. I joined an established Teamcenter X codebase and delivery workflow with no prior context, and my first weeks were spent reading other people's code and asking questions until the moving parts made sense. The team was upgrading production environments by hand, one at a time, guided by brittle legacy Ruby scripts, and the process was slow and error prone. As I got more comfortable, I started rewriting that automation in Python, replacing the Ruby scripts and cutting deployment time by 40 percent. The harder problem showed up during the upgrades themselves: I ran 44 production environment upgrades in four months, and each one surfaced live issues across Linux, databases, and services that I had to diagnose under time pressure, often by reproducing the failure locally first. That experience pushed me to design an Ansible based Blue-Green deployment framework that cut upgrade downtime in half and gave us a reliable way to roll back when something went wrong. Along the way I served as a quality coach for the team, raised automated test coverage from 35 to 78 percent, and mentored five engineers, so I learned to explain technical decisions clearly and take feedback well.

Please find my CV attached below. I look forward to hearing from you. I would be glad to contribute to the engineering behind the game and the VR planning tool and to grow into greater ownership as I learn the project.

{\setlength{\parskip}{0pt}
Regards,\\
Tanmay Bhuskute\\
bhuskutv@mail.uc.edu $\mid$ tanmay.v.bhuskute@gmail.com $\mid$ +1 (732) 858-3027\\
3424 Brookline Ave, Flat 19, Cincinnati, OH 45220
}

\end{letter}

\end{document}
